\usepackage{graphicx} % Required for inserting images
\usepackage{svg}
\usepackage{amsopn}
\usepackage{hyperref}
\usepackage{amsmath}
\usepackage{amssymb}
\usepackage{amstext}
\usepackage{amsfonts}
\usepackage{amsthm}
\usepackage{mathtools}
\usepackage{bbm}
\usepackage{makecell}
\usepackage{booktabs, multirow}
\usepackage{siunitx}
\usepackage{enumitem}
\usepackage{array}
\usepackage{varwidth}
\usepackage{todonotes}
\usepackage{floatrow}
\usepackage{bm}
\usepackage{xspace}
\usepackage{subcaption}
\usepackage[most]{tcolorbox}
\usepackage[table]{xcolor}
\usepackage{changepage}
\usepackage{algorithm}
\usepackage{algpseudocode}
\usepackage{pdflscape}
\usepackage{parskip}

\usepackage{mathpazo}
\usepackage{microtype}
\usepackage{dsfont}
%\usepackage[twoside]{geometry}
%\usepackage[twoside, bindingoffset=12mm]{geometry}

\newlength{\origtw}\setlength{\origtw}{\textwidth}
\newlength{\origth}\setlength{\origth}{\textheight}

\usepackage[twoside,
            textwidth=\origtw, textheight=\origth,
            hcentering,          % equal inner/outer margins
            bindingoffset=6mm    % per page, added at the spine
           ]{geometry}

\setlength{\parskip}{1em}
\setlength{\parindent}{1.5em}

% Create a shortcut for your "New Content" block
\newtcolorbox{newcontent}{
    colback=green!10!white, % 10% green, 90% white
    colframe=white,         % Hide the border
    sharp corners,          % Make it look like a highlight, not a box
    breakable,              % Let it span multiple pages
    enhanced,
    after skip=1em,         % Space after the block
    before skip=1em         % Space before the block
}

\sisetup{
  group-separator = {\,},
  group-minimum-digits = 4 % Ensures it separates even 4-`digit numbers like 1 800
}
\newcolumntype{V}[1]{>{\begin{varwidth}{#1}\centering\arraybackslash}c<{\end{varwidth}}}

\floatsetup[table]{objectset=centering}

\floatsetup[table]{objectset=centering, capposition=top}
\floatsetup[figure]{capposition=bottom}


%% Allow the S column to handle the "(std)" part without crashing
%\usepackage{siunitx}
%
%\sisetup{
%  separate-uncertainty=true,
%  table-align-uncertainty=true,
%  table-format = 1.2(4)
%}
\newcommand{\X}{\mathbf{X}}
\newcommand{\K}{\mathcal{K}}
\newcommand{\J}{\mathcal{J}}
\newcommand{\cL}{\mathcal{L}}
\renewcommand{\L}{\cL}
\newcommand{\I}{\mathcal{I}}
\newcommand{\Ind}{\mathds{1}}
\newcommand{\cP}{\mathcal{P}}
\newcommand{\cN}{\mathcal{N}}
\newcommand{\PN}{\mathcal{P} \cup \cN}
\renewcommand{\P}{\cP}
\newcommand{\N}{\cN}
\newcommand{\pun}{\P \cup \N}
\newcommand{\bs}{\mathcal{S}}
\newcommand{\R}{\mathbb{R}}
\newcommand{\Z}{\mathbb{Z}}


\renewcommand{\min}{\text{\textbf{min }}}
\renewcommand{\max}{\text{\textbf{max }}}
\newcommand{\din}{d^{\text{in}}}
\newcommand{\dout}{d^{\text{out}}}
\newcommand{\opt}{\text{OPT}}
\newcommand{\zopt}[1][]{\ifx&#1&\bm{\overline{z}}_{\opt}\else{\overline{z}}_{\opt,#1}\fi}
\newcommand{\gopt}[1][]{\ifx&#1&\bm{\overline{\gamma}}_{\opt}\else{\overline{\gamma}}_{\opt,#1}\fi}
\newcommand{\sopt}{S_\opt}
\newcommand{\st}{\text{\textbf{s.t.}}\;}
\newcommand{\zo}{0\mbox{-}1\xspace}
\newcommand{\fw}{\,:\,}
\newcommand{\topt}{\tau^{\opt}}
\newcommand{\tg}{\tau^{g}}

\newcommand{\umin}[1]{\mathop{\text{\textbf{min}}}\limits_{\,#1}}
\newcommand{\umax}[1]{\mathop{\text{\textbf{max}}}\limits_{\,#1}}
\newcommand{\smin}[1]{\text{\textbf{max}}_{#1}}
\newcommand{\smax}[1]{\text{\textbf{max}}_{#1}}
\newcommand{\conv}{\text{conv}}

\newcommand{\nq}{n_{\text{quantiles}}}

\newtheoremstyle{mydef}
  {1em}{1em}   % space above, below (independent of \partopsep)
  {}           % body font
  {}           % indent
  {\bfseries}  % head font
  {.}          % punctuation after head
  {.5em}       % space after head
  {}           % head spec
\theoremstyle{mydef}
\newtheorem{definition}{Definition}[section]
\newtheorem{claim}{Claim}[section]
\newtheorem{example}{Example}[section]
\newtheorem{lemma}{Lemma}[section]
\newtheorem{proposition}{Proposition}[section]
\newtheorem{theorem}{Theorem}[section]
\newtheorem{observation}{Observation}[section]
\newtheorem{conjecture}{Conjecture}[section]

%\renewcommand{\tableautorefname}{Table}
\renewcommand{\sectionautorefname}{Section}
\renewcommand{\subsectionautorefname}{Sub-section}
\renewcommand{\chapterautorefname}{Chapter}
\newcommand{\exampleautorefname}{Example}
\newcommand{\claimautorefname}{Claim}
\newcommand{\lemmaautorefname}{Lemma}
\newcommand{\propositionautorefname}{Proposition}
%\newcommand{\theoremautorefname}{Proposition}
\newcommand{\observationautorefname}{Observation}
\newcommand{\conjectureautorefname}{Conjecture}

\newcommand{\tquad}{\quad \quad \quad}

\DeclarePairedDelimiter{\abs}{\lvert}{\rvert}

\usepackage[backend=biber, style=numeric]{biblatex}
\addbibresource{references.bib}
\AtEveryBibitem{%
  \ifentrytype{article}{%
    \clearfield{url}
    \clearfield{urldate}
  }{}
}

\newcommand{\bwtable}[1]{%
    \begin{adjustwidth}{-2cm}{-2cm}
        \centering
        \begingroup
        \rowcolors{2}{gray!10}{white}
        \arrayrulecolor{black}
        \input{#1}
        \endgroup
    \end{adjustwidth}%
}

\newenvironment{note}{\par\noindent\textbf{Note: }}{\par}

%\usepackage{noindentafter}
%\NoIndentAfterEnv{definition}


